\chapter{Discussion}
\label{ch:discussion}

This chapter discusses the implications of our findings, addresses the research questions, and acknowledges the limitations of this work.

\section{Summary of Findings}

\subsection{LLM Model Selection}

Our evaluation of five large language models revealed that Mistral 7B achieved the highest performance for policy rule identification with 99.0\% accuracy and 0.98 confidence score (Table \ref{tab:model-summary}).

\begin{table}[h]
\centering
\caption{Summary of LLM Model Performance}
\label{tab:model-summary}
\begin{tabular}{|l|c|c|c|}
\hline
\textbf{Model} & \textbf{Accuracy} & \textbf{Confidence} & \textbf{Errors} \\
\hline
\textbf{Mistral 7B} & \textbf{99.0\%} & \textbf{0.98} & \textbf{0} \\
Mixtral 47B & 97.9\% & 0.91 & 2 \\
Llama 3.2 3B & 95.9\% & 0.88 & 4 \\
Phi-3 3.8B & 94.8\% & 0.91 & 4 \\
Llama 3.1 70B & 92.8\% & 0.89 & 2 \\
\hline
\end{tabular}
\end{table}

A notable finding is that Mistral 7B outperformed Llama 3.1 70B despite having 10x fewer parameters. This suggests that for specialized classification tasks, instruction-tuned smaller models may be more effective than larger general-purpose models.

\subsection{FOL Formalization Success}

The FOL formalization achieved 100\% success rate after technical optimizations (v2), up from 91.7\% in v1 (Table \ref{tab:fol-versions}).

\begin{table}[h]
\centering
\caption{FOL Formalization Version Comparison}
\label{tab:fol-versions}
\begin{tabular}{|l|c|c|c|}
\hline
\textbf{Version} & \textbf{Success} & \textbf{Failed} & \textbf{Rate} \\
\hline
v1 (Original) & 88 & 8 & 91.7\% \\
\textbf{v2 (Improved)} & \textbf{96} & \textbf{0} & \textbf{100\%} \\
\hline
\end{tabular}
\end{table}

The 8 initially failed cases were recovered through improved JSON parsing and increased response length, not through changes to the logic representation.

\section{Research Questions Addressed}

\subsection{RQ1: Can LLMs Identify Policy Rules?}

\textbf{Answer: Yes, with high accuracy.}

Our results demonstrate that LLMs can effectively identify policy rules in academic documents. Mistral 7B correctly classified 96 out of 97 rules (99\%), with inter-model agreement rates exceeding 88\% across all model pairs.

\subsection{RQ2: Is First-Order Logic Sufficient for Policy Formalization?}

\textbf{Answer: Yes, FOL with deontic operators is sufficient.}

Our analysis found that 100\% of policy rules could be formalized using FOL extended with deontic operators (O, P, F). We evaluated whether Second-Order Logic (SOL) or Higher-Order Logic (HOL) would be necessary and concluded they are not:

\begin{itemize}
    \item The initial 8.3\% failure rate was due to JSON parsing issues, not expressiveness limitations
    \item No rules required quantification over predicates (SOL) or functions as objects (HOL)
    \item Temporal and numerical extensions within FOL are sufficient for deadline and threshold constraints
\end{itemize}

This finding aligns with Governatori \& Rotolo (2010), who state that ``first-order deontic logic provides sufficient expressiveness for normative systems.''

\subsection{RQ3: Can Policy Rules be Translated to SHACL?}

\textbf{Answer: Yes, with systematic translation.}

We successfully translated 96 FOL-formalized rules to SHACL shapes:
\begin{itemize}
    \item 65 obligations (68\%)
    \item 17 permissions (18\%)
    \item 14 prohibitions (15\%)
\end{itemize}

The translation preserves deontic semantics through SHACL's severity levels and custom deontic type annotations.

\section{Why SOL and HOL Are Not Required}

\subsection{Second-Order Logic (SOL)}

SOL allows quantification over predicates and sets. Our analysis found no policy rules requiring such constructs:

\begin{itemize}
    \item All subject classes are concrete (Student, Employee, Sponsor)
    \item All conditions are first-order (enrolled, paid, approved)
    \item No meta-level reasoning about predicates themselves
\end{itemize}

Furthermore, SOL introduces undecidability (Brachman \& Levesque, 2004), making automated verification impractical.

\subsection{Higher-Order Logic (HOL)}

HOL treats functions as first-class objects. This is unnecessary because:

\begin{itemize}
    \item Policy rules reference concrete actions, not abstract functions
    \item No lambda expressions or function composition needed
    \item Type polymorphism is not required for institutional policies
\end{itemize}

As Farmer (2008) notes, HOL's expressiveness is designed for mathematical theorem proving, not normative reasoning.

\subsection{What Extensions Are Needed}

Instead of SOL/HOL, we recommend domain-specific extensions:

\begin{table}[h]
\centering
\caption{Recommended FOL Extensions}
\label{tab:extensions}
\begin{tabular}{|l|l|l|}
\hline
\textbf{Extension} & \textbf{Purpose} & \textbf{Example} \\
\hline
Temporal Logic & Deadlines & Within(2, weeks, vacate) \\
Constraint Logic & Thresholds & value >= 3000 \\
Defeasible Logic & Exceptions & unless(approved) \\
\hline
\end{tabular}
\end{table}

\section{Limitations}

\subsection{Technical Limitations}

\begin{enumerate}
    \item \textbf{Single Institution}: Tested only on AIT policy documents
    \item \textbf{English Only}: No evaluation of multilingual policies
    \item \textbf{Manual Parsing Recovery}: 8 rules required fallback parsing
\end{enumerate}

\subsection{Methodological Limitations}

\begin{enumerate}
    \item \textbf{No Human Gold Standard}: LLM agreement used as proxy for accuracy
    \item \textbf{Limited Validation}: SHACL shapes not tested against real institutional data
    \item \textbf{Static Analysis}: No runtime policy enforcement evaluated
\end{enumerate}

\section{Implications for Practice}

\subsection{For Policy Administrators}

\begin{itemize}
    \item LLMs can assist in identifying and extracting policy rules
    \item Automated formalization is feasible with appropriate tooling
    \item SHACL-based validation can check compliance programmatically
\end{itemize}

\subsection{For Researchers}

\begin{itemize}
    \item Smaller instruction-tuned models may outperform larger ones for classification
    \item FOL remains adequate for institutional policy formalization
    \item Practical engineering issues (parsing, response length) may be more critical than logic expressiveness
\end{itemize}

\section{Future Work}

\begin{enumerate}
    \item \textbf{Multi-institutional Validation}: Test on policies from different universities
    \item \textbf{Human Annotation Study}: Create gold standard with expert annotators
    \item \textbf{Runtime Enforcement}: Integrate SHACL validation with student information systems
    \item \textbf{Policy Conflict Detection}: Extend to identify contradictory rules
    \item \textbf{Multilingual Support}: Evaluate on non-English policy documents
\end{enumerate}
